\uuid{5dK8}
\titre{ Loi à densité }

\niveau{L3}
\module{ Probabilités }
\chapitre{Lois absolument continues}
\sousChapitre{}
\theme{}
\auteur{Quentin Liard}
\datecreate{ 2026-04-06}
\organisation{AMSCC}
\difficulte{2}

\contenu{

\texte{ 
On considère la fonction $f$ définie sur $\mathbb{R}$ par
\[
f(x)=
\begin{cases}
c\,x(3-x) & \text{si } x\in[0,2],\\
0 & \text{sinon},
\end{cases}
\]
où $c$ est une constante réelle. On suppose que $f$ est la densité d’une variable aléatoire réelle $X$.
}

\begin{enumerate}

\item   \question{Déterminer la constante $c$ pour que $f$ soit bien une densité de probabilité sur $\mathbb{R}$.}
\indication{}
\reponse{
Pour que $f$ soit une densité, il faut
\[
\int_{-\infty}^{+\infty} f(x)\,dx=1.
\]
Comme $f(x)=0$ hors de $[0,2]$, on doit avoir
\[
\int_0^2 c\,x(3-x)\,dx=1.
\]
Or
\[
x(3-x)=3x-x^2,
\]
donc
\[
\int_0^2 x(3-x)\,dx=\int_0^2 (3x-x^2)\,dx
=\left[\frac{3x^2}{2}-\frac{x^3}{3}\right]_0^2
=6-\frac{8}{3}
=\frac{10}{3}.
\]
Ainsi
\[
c\cdot \frac{10}{3}=1,
\]
d’où
\[
c=\frac{3}{10}.
\]
}

%\item   \question{Calculer la fonction de répartition $F_X(x)=P(X\le x)$.}
%\indication{}
%\reponse{
%On a
%\[
%F_X(x)=\int_{-\infty}^x f(t)\,dt.
%\]
%On distingue les cas :

%Si $x<0$, alors
%\[
%F_X(x)=0.
%\]

%Si $0\le x\le 2$, alors
%\[
%F_X(x)=\int_0^x \frac{3}{10}\,t(3-t)\,dt
%=\frac{3}{10}\int_0^x (3t-t^2)\,dt.
%\]
%Donc
%\[
%F_X(x)=\frac{3}{10}\left[\frac{3t^2}{2}-\frac{t^3}{3}\right]_0^x
%=\frac{3}{10}\left(\frac{3x^2}{2}-\frac{x^3}{3}\right).
%\]
%Ainsi
%\[
%F_X(x)=\frac{9x^2}{20}-\frac{x^3}{10}.
%\]

%Si $x>2$, alors
%\[
%F_X(x)=1.
%\]

%Finalement,
%\[
%F_X(x)=
%\begin{cases}
%0 & \text{si } x<0,\\[0.4em]
%\displaystyle \frac{9x^2}{20}-\frac{x^3}{10} & \text{si } 0\le x\le 2,\\[0.8em]
%1 & \text{si } x>2.
%\end{cases}
%\]
%}

\item   \question{Calculer $P(1\le X\le 2)$.}
\indication{}
\reponse{
\[
P(1\le X\le 2)=\int_1^2 \frac{3}{10}\,x(3-x)\,dx.
\]
On peut aussi utiliser la fonction de répartition :
\[
P(1\le X\le 2)=F_X(2)-F_X(1).
\]
Or
\[
F_X(2)=1,
\qquad
F_X(1)=\frac{9}{20}-\frac{1}{10}=\frac{7}{20}.
\]
Donc
\[
P(1\le X\le 2)=1-\frac{7}{20}=\frac{13}{20}.
\]
}

\item   \question{Calculer l’espérance $E(X)$ et $E(X^2)$.}
\indication{}
\reponse{
On a
\[
E(X)=\int_{-\infty}^{+\infty} x f(x)\,dx
=\int_0^2 x\cdot \frac{3}{10}x(3-x)\,dx
=\frac{3}{10}\int_0^2 x^2(3-x)\,dx.
\]
Donc
\[
E(X)=\frac{3}{10}\int_0^2 (3x^2-x^3)\,dx
=\frac{3}{10}\left[x^3-\frac{x^4}{4}\right]_0^2
=\frac{3}{10}(8-4)
=\frac{6}{5}.
\]

De même,
\[
E(X^2)=\int_{-\infty}^{+\infty} x^2 f(x)\,dx
=\int_0^2 x^2\cdot \frac{3}{10}x(3-x)\,dx
=\frac{3}{10}\int_0^2 x^3(3-x)\,dx.
\]
Ainsi
\[
E(X^2)=\frac{3}{10}\int_0^2 (3x^3-x^4)\,dx
=\frac{3}{10}\left[\frac{3x^4}{4}-\frac{x^5}{5}\right]_0^2.
\]
Or
\[
\frac{3\cdot 2^4}{4}-\frac{2^5}{5}=12-\frac{32}{5}=\frac{28}{5}.
\]
Donc
\[
E(X^2)=\frac{3}{10}\cdot \frac{28}{5}=\frac{42}{25}.
\]
}

\item   \question{En déduire la variance $V(X)$.}
\indication{}
\reponse{
On utilise la formule
\[
V(X)=E(X^2)-[E(X)]^2.
\]
Donc
\[
V(X)=\frac{42}{25}-\left(\frac{6}{5}\right)^2
=\frac{42}{25}-\frac{36}{25}
=\frac{6}{25}.
\]
Ainsi
\[
V(X)=\frac{6}{25}.
\]
}

\end{enumerate}

}